\subsection{What is an agent architecture?}
\label{sec:architectures-overview}
\slidesrc{Agent architectures: 11--14}

\begin{nbdef}{Agent architecture (slide 11)}
An \textbf{architecture} proposes a particular methodology for building an autonomous agent:
\begin{itemize}[nosep]
  \item how the construction of the agent can be decomposed into the construction of a set
        of \textbf{component modules};
  \item how these modules should be made to \textbf{interact}.
\end{itemize}
These two aspects define how the \emph{sensor data} and the \emph{current internal state}
of the agent determine the \emph{actions} (effector outputs) and the \emph{future internal
state} of the agent.
\end{nbdef}

\begin{center}
\begin{tikzpicture}[font=\small, >=Stealth,
  fn/.style={draw, thick, circle, minimum size=0.9cm, fill=nbblue!10},
  st/.style={draw, thick, minimum width=0.8cm, minimum height=0.9cm, fill=nborange!12}]
  \draw[thick, nbgray, rounded corners] (0.6,-1.4) rectangle (7.4,1.5);
  \node[font=\scriptsize, nbgray] at (6.9,-1.2) {agent};
  \node (in) at (-0.8,0) {perception $p$};
  \node[fn] (f) at (2,0) {$f$};
  \node[st] (s) at (4,0) {$s$};
  \node[fn] (g) at (6,0) {$g$};
  \node (out) at (8.8,0) {action $a$};
  \draw[->, thick] (in) -- (f);
  \draw[->, thick] (f) -- (s);
  \draw[->, thick] (s) -- (g);
  \draw[->, thick] (g) -- (out);
  \draw[->, thick] (1.0,0) |- (6,1.05) -- (g.north);
  \draw[->, thick] (4.7,0) |- (2,-0.95) -- (f.south);
  \fill (1.0,0) circle (1.6pt); \fill (4.7,0) circle (1.6pt);
\end{tikzpicture}
\captionof{figure}{From perception to action (slide 12): $f$ = state update function, $s$ =
internal state, $g$ = output function. The new state depends on the old state and the
percept; the action depends on the state and the percept.}
\end{center}

In formulas, at each step $t$:
\begin{align*}
s_t &= f(s_{t-1},\, p_t) && \reason{update the internal state with what was perceived} \\
a_t &= g(s_t,\, p_t) && \reason{choose the action}
\end{align*}

\begin{nbexample}{a smart thermostat as $f$, $s$, $g$}
\notinslides
Percept $p_t$: the measured temperature. State $s_t$: the average of the last three
temperatures (to ignore a sensor glitch). Output: switch the heater on if the average is
below 20\,$^\circ$C.
\begin{center}
\small
\begin{tabular}{@{}c c c c l@{}}
\toprule
$t$ & $p_t$ & last three & $s_t$ (average) & $a_t = g(s_t, p_t)$ \\
\midrule
1 & 19 & 21, 21, 19 & 20.33 & off \\
2 & 18 & 21, 19, 18 & 19.33 & on \\
3 & 25 (glitch) & 19, 18, 25 & 20.67 & off \\
\bottomrule
\end{tabular}
\end{center}
(Initial readings 21, 21.) The internal state is what lets an agent do better than reacting
to the last percept only; how much state, and how it is used, is exactly what distinguishes
the architectures below.
\end{nbexample}

\begin{nbfact}{Main kinds of agent architectures (slide 13)}
\begin{itemize}[nosep]
  \item \textbf{Reactive} architectures: focused on fast reactions/responses to changes
        detected in the environment.
  \item \textbf{Deliberative} (symbolic) architectures: focused on long-term planning of
        actions, centred on a set of basic goals.
  \item \textbf{Hybrid} architectures: combining a reactive side and a deliberative side.
\end{itemize}
\end{nbfact}

\begin{nbexample}{reactive vs.\ deliberative: a robot that has to reach a point (slide 14)}
\begin{center}
\small
\begin{tabularx}{\linewidth}{@{}X X@{}}
\toprule
\textbf{Reactive robot} & \textbf{Deliberative robot} \\
\midrule
a sensor in the front of the robot & an explicit representation of the environment (a map) \\
change movement right/left when the sensor detects an obstacle & a planning procedure finds
  the minimal route between the current position and the destination \\
minimal computation, based on current location and destination point & high computational
  cost; possible dynamic re-planning needed \\
\bottomrule
\end{tabularx}
\end{center}
\begin{center}
\begin{tikzpicture}[scale=0.62, font=\scriptsize, >=Stealth]
  \begin{scope}
    \draw[nbgray] (0,0) rectangle (6,4);
    \fill[nbgray!50] (2,1) rectangle (3,4); \fill[nbgray!50] (4,0) rectangle (4.6,2.6);
    \node[circle, fill=nbblue, inner sep=2pt, label=below:start] (A) at (0.6,0.6) {};
    \node[star, fill=nbgreen, inner sep=1.5pt, label=above:goal] (B) at (5.4,3.4) {};
    \draw[->, thick, nborange] (A) -- (1.8,0.6) -- (1.8,0.4) -- (2.2,0.5) -- (3.4,0.4) -- (3.6,1.4) -- (3.8,2.9) -- (4.4,2.8) -- (4.9,3.2) -- (B);
    \node at (3,-0.6) {reactive: bumps and turns};
  \end{scope}
  \begin{scope}[xshift=8cm]
    \draw[nbgray] (0,0) rectangle (6,4);
    \fill[nbgray!50] (2,1) rectangle (3,4); \fill[nbgray!50] (4,0) rectangle (4.6,2.6);
    \node[circle, fill=nbblue, inner sep=2pt, label=below:start] (A) at (0.6,0.6) {};
    \node[star, fill=nbgreen, inner sep=1.5pt, label=above:goal] (B) at (5.4,3.4) {};
    \draw[->, thick, nbpurple] (A) -- (3.5,0.5) -- (3.6,3.1) -- (B);
    \node at (3,-0.6) {deliberative: planned shortest path};
  \end{scope}
\end{tikzpicture}
\captionof{figure}{Same task, two architectures. The reactive robot computes almost nothing
but wanders; the deliberative one finds the best route but must have a correct map and
re-plan if a wall moves.}
\end{center}
\end{nbexample}
